\paragrafos{Άθροισμα γωνιών τριγώνου}{Β.3.9, «Άθροισμα γωνιών τριγώνου»}

\begin{anagke}
Η κυρία Αλεξάνδρα ζωγραφίζει στον πίνακα ένα τρίγωνο και ζητά από τον Στέλιο να μετρήσει τις γωνίες του. Εκείνος μετρά δύο γωνίες, $62^\circ$ και $48^\circ$, αλλά το μοιρογνωμόνιο γλιστράει και η τρίτη μέτρηση χάνεται. «Δεν πειράζει», λέει η Φανή, «δεν χρειάζεται να τη μετρήσουμε ξανά». Ο Στέλιος δεν πείθεται: «Και πώς θα τη βρούμε χωρίς να τη μετρήσουμε;». Η Ζωή ρωτά: «Μήπως οι τρεις γωνίες ενός τριγώνου έχουν πάντα το ίδιο άθροισμα, όποιο τρίγωνο κι αν σχεδιάσεις;». Η κυρία Αλεξάνδρα χαμογελά: «Κόψτε ένα τρίγωνο από χαρτί, σκίστε τις τρεις γωνίες και βάλτε τες δίπλα δίπλα με τις κορυφές στο ίδιο σημείο. Τι σχηματίζεται;».

Σε αυτή την παράγραφο μαθαίνουμε ότι το άθροισμα των γωνιών κάθε τριγώνου είναι πάντα το ίδιο, και το χρησιμοποιούμε για να βρίσκουμε γωνίες που δεν μπορούμε ή δεν θέλουμε να μετρήσουμε.
\end{anagke}

\begin{lexi}{άθροισμα, προσκείμενες γωνίες, συμπληρωματικές}{άθροισμα γωνιών}
Το \lexiword{άθροισμα} είναι το «ύψος» που φτάνει ένας σωρός όταν βάζεις τα πράγματα ένα πάνω στο άλλο, δηλαδή το αποτέλεσμα της πρόσθεσης. Οι \lexiword{προσκείμενες γωνίες} μιας πλευράς είναι όσες «κάθονται» πάνω της (έχουν την πλευρά αυτή κοινή). Οι \lexiword{συμπληρωματικές} γωνίες «συμπληρώνουν» η μία την άλλη μέχρι την ορθή γωνία.
\end{lexi}

\begin{metafrasi}
Οι τρεις γωνίες του τριγώνου μαζί. & $\hat{\mathrm{Α}}+\hat{\mathrm{Β}}+\hat{\mathrm{Γ}}$\\
Μια ευθεία γωνία έχει μέτρο. & $180^\circ$\\
Οι γωνίες που κάθονται στη βάση του ισοσκελούς. & προσκείμενες στη βάση, $\hat{\mathrm{Β}}=\hat{\mathrm{Γ}}$\\
Δύο γωνίες που μαζί κάνουν ορθή. & συμπληρωματικές ($90^\circ$)\\
\end{metafrasi}

\tmima{Η θεωρία}

Το πείραμα με το χαρτί δείχνει ότι οι τρεις γωνίες ενός τριγώνου σχηματίζουν μαζί μια ευθεία γωνία. Μπορούμε να το αποδείξουμε με όσα ξέρουμε για τις παράλληλες ευθείες.

\begin{kanonas}{άθροισμα γωνιών τριγώνου}
Το άθροισμα των γωνιών κάθε τριγώνου είναι $180^\circ$:
\[\hat{\mathrm{Α}}+\hat{\mathrm{Β}}+\hat{\mathrm{Γ}}=180^\circ.\]
\end{kanonas}

\begin{center}
\begin{tikzpicture}[scale=0.95]
\draw[thick] (0,0)--(5,0)--(1.6,2.4)--cycle;
\draw[gblue,thick] (-0.8,2.4)--(5.8,2.4) node[right] {$\varepsilon$};
\node[below left] at (0,0) {$\mathrm{Β}$}; \node[below right] at (5,0) {$\mathrm{Γ}$}; \node[above] at (1.6,2.4) {$\mathrm{Α}$};
\fill (1.6,2.4) circle (1.6pt);
\node at (0.8,0.2) {\small $\hat{\mathrm{Β}}$};
\node at (4.2,0.2) {\small $\hat{\mathrm{Γ}}$};
\node at (1.5,1.9) {\small $\hat{\mathrm{Α}}$};
\node at (0.85,2.15) {\small $\hat{\mathrm{Α}}_1$};
\node at (2.9,2.15) {\small $\hat{\mathrm{Α}}_2$};
\end{tikzpicture}
\end{center}

\textbf{Απόδειξη.} Από την κορυφή $\mathrm{Α}$ φέρνουμε την ευθεία $\varepsilon$ παράλληλη προς τη $\mathrm{ΒΓ}$. Τότε $\hat{\mathrm{Α}}_1=\hat{\mathrm{Β}}$ (εντός εναλλάξ, με τέμνουσα την $\mathrm{ΑΒ}$) και $\hat{\mathrm{Α}}_2=\hat{\mathrm{Γ}}$ (εντός εναλλάξ, με τέμνουσα την $\mathrm{ΑΓ}$). Οι γωνίες $\hat{\mathrm{Α}}_1$, $\hat{\mathrm{Α}}$, $\hat{\mathrm{Α}}_2$ σχηματίζουν μαζί ευθεία γωνία, άρα $\hat{\mathrm{Α}}_1+\hat{\mathrm{Α}}+\hat{\mathrm{Α}}_2=180^\circ$. Αντικαθιστώντας, παίρνουμε $\hat{\mathrm{Β}}+\hat{\mathrm{Α}}+\hat{\mathrm{Γ}}=180^\circ$.

Από τον κανόνα προκύπτουν χρήσιμα συμπεράσματα.

\begin{idiotita}{συνέπειες}
\begin{itemize}[leftmargin=1.2em,itemsep=0.3ex,topsep=2pt]
\item Σε \textbf{ισόπλευρο} τρίγωνο όλες οι γωνίες είναι ίσες, άρα κάθε γωνία είναι $180^\circ:3=60^\circ$.
\item Σε \textbf{ορθογώνιο} τρίγωνο οι δύο οξείες γωνίες είναι \textbf{συμπληρωματικές}: $\hat{\mathrm{Β}}+\hat{\mathrm{Γ}}=90^\circ$.
\item Σε \textbf{ορθογώνιο και ισοσκελές} τρίγωνο οι οξείες γωνίες είναι $45^\circ$ και $45^\circ$.
\item Σε \textbf{ισοσκελές} τρίγωνο οι γωνίες της βάσης είναι ίσες, άρα
\[\text{κορυφή}=180^\circ-2\cdot(\text{βάση}),\]
\[\text{βάση}=\bigl(180^\circ-\text{κορυφή}\bigr):2.\]
\item Ένα τρίγωνο έχει \textbf{το πολύ μία} ορθή ή αμβλεία γωνία, και \textbf{τουλάχιστον δύο} οξείες.
\end{itemize}
\end{idiotita}

\begin{prosoxi}
\begin{itemize}[leftmargin=1.2em,itemsep=0.3ex,topsep=2pt]
\item Το άθροισμα είναι πάντα $180^\circ$, όσο μεγάλο ή μικρό κι αν είναι το τρίγωνο. Το μέγεθος δεν αλλάζει το άθροισμα.
\item Αν σε ισοσκελές τρίγωνο σου δίνουν μία γωνία, πρέπει να σκεφτείς αν είναι η γωνία της κορυφής ή της βάσης. Συχνά υπάρχουν \emph{δύο} δυνατές περιπτώσεις.
\item Το ισοσκελές τρίγωνο έχει ίσες μόνο τις γωνίες της βάσης. Και οι τρεις γωνίες είναι ίσες μόνο στο ισόπλευρο.
\end{itemize}
\end{prosoxi}

\tmima{Λυμένα παραδείγματα}

\paradeigma{Η τρίτη γωνία του Στέλιου}
\tixero Σε τρίγωνο δύο γωνίες είναι $62^\circ$ και $48^\circ$.\par
\tithelo Η τρίτη γωνία και το είδος του τριγώνου.\par
\possindeo Το άθροισμα των τριών γωνιών είναι $180^\circ$. Οι δύο γωνίες μαζί κάνουν $62^\circ+48^\circ=110^\circ$, άρα η τρίτη είναι $180^\circ-110^\circ=70^\circ$.\par
\apantisi $70^\circ$. Όλες οι γωνίες είναι οξείες, άρα το τρίγωνο είναι οξυγώνιο.\par
\elegxos $62^\circ+48^\circ+70^\circ=180^\circ$.

\paradeigma{Ισοσκελές με γνωστή κορυφή}
\tixero Ισοσκελές τρίγωνο $\mathrm{ΑΒΓ}$ ($\mathrm{ΑΒ}=\mathrm{ΑΓ}$) με $\hat{\mathrm{Α}}=50^\circ$.\par
\tithelo Οι γωνίες $\hat{\mathrm{Β}}$ και $\hat{\mathrm{Γ}}$.\par
\possindeo Οι γωνίες της βάσης είναι ίσες, $\hat{\mathrm{Β}}=\hat{\mathrm{Γ}}=x$. Τότε $50^\circ+x+x=180^\circ$, άρα $2x=130^\circ$ και $x=65^\circ$.\par
\apantisi $\hat{\mathrm{Β}}=\hat{\mathrm{Γ}}=65^\circ$.\par
\elegxos $50^\circ+65^\circ+65^\circ=180^\circ$.

\paradeigma{Ισοσκελές με μία γνωστή γωνία}
\tixero Ένα ισοσκελές τρίγωνο έχει μια γωνία $80^\circ$.\par
\tithelo Οι άλλες δύο γωνίες.\par
\possindeo Η γωνία των $80^\circ$ μπορεί να είναι η γωνία της κορυφής ή μία γωνία της βάσης.\par
\textbf{1η περίπτωση: κορυφή.} Οι γωνίες της βάσης είναι ίσες: $(180^\circ-80^\circ):2=50^\circ$ η καθεμία.\par
\textbf{2η περίπτωση: βάση.} Η άλλη γωνία της βάσης είναι επίσης $80^\circ$, και η κορυφή είναι $180^\circ-2\cdot80^\circ=20^\circ$.\par
\apantisi $(80^\circ,50^\circ,50^\circ)$ ή $(80^\circ,80^\circ,20^\circ)$.\par
\elegxos Και στις δύο περιπτώσεις το άθροισμα είναι $180^\circ$.

\paradeigma{Γωνίες με άγνωστο}
\tixero Οι γωνίες ενός τριγώνου είναι $x$, $x+20^\circ$ και $2x-20^\circ$.\par
\tithelo Το $x$ και οι γωνίες του τριγώνου.\par
\possindeo Το άθροισμα είναι $180^\circ$: $x+(x+20^\circ)+(2x-20^\circ)=180^\circ$. Στο άθροισμα, τα $+20^\circ$ και $-20^\circ$ αλληλοαναιρούνται και μένει $4x=180^\circ$, άρα $x=45^\circ$.\par
\apantisi Οι γωνίες είναι $45^\circ$, $65^\circ$ και $70^\circ$.\par
\elegxos $45^\circ+65^\circ+70^\circ=180^\circ$. Το τρίγωνο είναι οξυγώνιο.

\paradeigma{Παράλληλες και τρίγωνα}
\tixero Οι $\mathrm{ΑΒ}$ και $\mathrm{ΓΔ}$ είναι παράλληλες και τα τμήματα $\mathrm{ΑΔ}$ και $\mathrm{ΒΓ}$ τέμνονται στο $\mathrm{Ο}$. Είναι $\widehat{\mathrm{ΟΑΒ}}=40^\circ$ και $\widehat{\mathrm{ΟΓΔ}}=35^\circ$.\par
\tithelo Η γωνία $\widehat{\mathrm{ΑΟΒ}}$.\par
\possindeo Η $\widehat{\mathrm{ΟΒΑ}}$ είναι εντός εναλλάξ με τη $\widehat{\mathrm{ΟΓΔ}}$ (παράλληλες $\mathrm{ΑΒ}$, $\mathrm{ΓΔ}$ και τέμνουσα $\mathrm{ΒΓ}$), άρα $\widehat{\mathrm{ΟΒΑ}}=35^\circ$. Στο τρίγωνο $\mathrm{ΟΑΒ}$ είναι $40^\circ+35^\circ+\widehat{\mathrm{ΑΟΒ}}=180^\circ$, άρα $\widehat{\mathrm{ΑΟΒ}}=105^\circ$.\par
\apantisi $\widehat{\mathrm{ΑΟΒ}}=105^\circ$.\par
\elegxos Η $\widehat{\mathrm{ΓΟΔ}}$ είναι κατακορυφήν με την $\widehat{\mathrm{ΑΟΒ}}$, άρα $105^\circ$, και στο τρίγωνο $\mathrm{ΟΓΔ}$ έχουμε $35^\circ+40^\circ+105^\circ=180^\circ$.

\begin{center}
\begin{tikzpicture}[scale=0.85]
\draw[thick] (-1.6,1.34)--(1.92,1.34);
\draw[thick] (-1.92,-1.34)--(1.6,-1.34);
\draw[thick,gblue] (-1.6,1.34)--(1.6,-1.34);
\draw[thick,gblue] (1.92,1.34)--(-1.92,-1.34);
\fill (-1.6,1.34) circle (1.6pt) node[above left] {$\mathrm{Α}$};
\fill (1.92,1.34) circle (1.6pt) node[above right] {$\mathrm{Β}$};
\fill (-1.92,-1.34) circle (1.6pt) node[below left] {$\mathrm{Γ}$};
\fill (1.6,-1.34) circle (1.6pt) node[below right] {$\mathrm{Δ}$};
\fill (0,0) circle (1.6pt) node[above] {\small $\mathrm{Ο}$};
\node at (-0.65,1.1) {\small $40^\circ$};
\node at (-1.1,-1.1) {\small $35^\circ$};
\end{tikzpicture}
\end{center}

\tmima{Ασκήσεις}

\zoni{Το κατάλαβα}{άθροισμα $180^\circ$}

\ex Σωστό ή λάθος; α) Υπάρχει τρίγωνο με δύο ορθές γωνίες. β) Υπάρχει τρίγωνο με γωνίες $50^\circ$, $60^\circ$ και $70^\circ$. γ) Οι οξείες γωνίες ενός ορθογωνίου τριγώνου είναι συμπληρωματικές. δ) Υπάρχει ισόπλευρο τρίγωνο με γωνία $70^\circ$. ε) Αν η γωνία της κορυφής ισοσκελούς τριγώνου είναι $30^\circ$, κάθε γωνία της βάσης είναι $75^\circ$. στ) Ένα τρίγωνο μπορεί να έχει δύο αμβλείες γωνίες.

\ex Βρες την τρίτη γωνία τριγώνου αν οι δύο άλλες είναι: α) $45^\circ$ και $65^\circ$ β) $90^\circ$ και $34^\circ$ γ) $108^\circ$ και $25^\circ$ δ) $59^\circ$ και $61^\circ$.

\ex Σε ισοσκελές τρίγωνο η γωνία της κορυφής είναι: α) $40^\circ$ β) $100^\circ$ γ) $90^\circ$ δ) $20^\circ$. Βρες τις γωνίες της βάσης.

\ex Σε ισοσκελές τρίγωνο η γωνία της βάσης είναι: α) $65^\circ$ β) $30^\circ$ γ) $45^\circ$ δ) $80^\circ$. Βρες τη γωνία της κορυφής.

\ex Οι γωνίες ενός τριγώνου είναι $x$, $x+10^\circ$ και $x+20^\circ$. Βρες το $x$ και τις τρεις γωνίες.

\ex Ποιες από τις παρακάτω τριάδες μπορεί να είναι οι γωνίες ενός τριγώνου; α) $70^\circ,60^\circ,50^\circ$ β) $90^\circ,50^\circ,50^\circ$ γ) $100^\circ,40^\circ,40^\circ$ δ) $120^\circ,30^\circ,40^\circ$ ε) $60^\circ,60^\circ,60^\circ$ στ) $35^\circ,45^\circ,90^\circ$.

\ex Πες αν το τρίγωνο με τις παρακάτω γωνίες είναι οξυγώνιο, ορθογώνιο ή αμβλυγώνιο: α) $50^\circ,60^\circ,70^\circ$ β) $90^\circ,35^\circ,55^\circ$ γ) $105^\circ,40^\circ,35^\circ$ δ) $80^\circ,50^\circ,50^\circ$.

\zoni{Το χρησιμοποιώ}{στέγες, σκάλες, χωράφια}

\ex \textbf{Η στέγη.} Το αέτωμα του εξοχικού της θείας Μαρίας έχει σχήμα ισοσκελούς τριγώνου με γωνία στην κορυφή $110^\circ$. Πόση γωνία σχηματίζει κάθε πλευρά της στέγης με το οριζόντιο δοκάρι της βάσης;

\ex \textbf{Το σημαιάκι.} Η Ρήνη κόβει ένα ισοσκελές τρίγωνο για σημαιάκι και μία γωνία του είναι $50^\circ$. Βρες όλες τις δυνατές τριάδες γωνιών που μπορεί να έχει.

\ex \textbf{Η σκάλα.} Ο Ηλίας ακουμπά μια σκάλα σε κατακόρυφο τοίχο. Η σκάλα σχηματίζει με το οριζόντιο δάπεδο γωνία $72^\circ$. Πόση γωνία σχηματίζει με τον τοίχο;

\ex \textbf{Το χωράφι.} Ένα τριγωνικό χωράφι του παππού Γιάννη έχει γωνίες $2x$, $3x$ και $4x$. Βρες τις γωνίες του και το είδος του τριγώνου ως προς τις γωνίες.

\ex \textbf{Το σχέδιο της Αγγελικής.} Η Αγγελική ζωγραφίζει ισοσκελές τρίγωνο όπου η γωνία της κορυφής είναι τριπλάσια από τη γωνία της βάσης. Βρες τις γωνίες του.

\zoni{Το σκέφτομαι}{γιατί και πώς}

\ex Εξήγησε γιατί κάθε τρίγωνο έχει τουλάχιστον δύο οξείες γωνίες.

\ex Το τρίγωνο $\mathrm{ΑΒΓ}$ είναι ισοσκελές με $\mathrm{ΑΒ}=\mathrm{ΑΓ}$ και $\hat{\mathrm{Α}}=44^\circ$. Η ευθεία $\mathrm{ΔΕ}$ είναι παράλληλη προς τη $\mathrm{ΒΓ}$, με το $\mathrm{Δ}$ στην $\mathrm{ΑΒ}$ και το $\mathrm{Ε}$ στην $\mathrm{ΑΓ}$. Βρες τις γωνίες $\widehat{\mathrm{ΑΔΕ}}$, $\widehat{\mathrm{ΑΕΔ}}$ και $\widehat{\mathrm{ΔΕΓ}}$.

\begin{center}
\begin{tikzpicture}[scale=0.62]
\draw[thick] (0,0)--(4,0)--(2,4.95)--cycle;
\draw[thick,gblue] (0.8,1.98)--(3.2,1.98);
\node[below left] at (0,0) {$\mathrm{Β}$}; \node[below right] at (4,0) {$\mathrm{Γ}$}; \node[above] at (2,4.95) {$\mathrm{Α}$};
\node[left] at (0.8,1.98) {$\mathrm{Δ}$}; \node[right] at (3.2,1.98) {$\mathrm{Ε}$};
\node at (2,4.2) {\small $44^\circ$};
\end{tikzpicture}
\end{center}

\ex Οι $\varepsilon_1$ και $\varepsilon_2$ είναι παράλληλες και μια τέμνουσα τις τέμνει στα $\mathrm{Α}$ και $\mathrm{Β}$. Οι εντός και επί τα αυτά γωνίες είναι $70^\circ$ στο $\mathrm{Α}$ και $110^\circ$ στο $\mathrm{Β}$. Οι διχοτόμοι τους τέμνονται στο $\mathrm{Κ}$. Βρες τη γωνία $\widehat{\mathrm{ΑΚΒ}}$. Θα έβγαινε το ίδιο με άλλες τιμές γωνιών;

\ex Στο τρίγωνο $\mathrm{ΑΒΓ}$ είναι $\hat{\mathrm{Β}}=30^\circ$ και η γωνία $\hat{\mathrm{Α}}$ είναι τετραπλάσια της $\hat{\mathrm{Γ}}$. Βρες τις γωνίες του τριγώνου και το είδος του ως προς τις γωνίες και τις πλευρές του.

\zoni{Για το σπίτι}{ρίξε μια ματιά}

\ex Κόψε από χαρτί ένα τρίγωνο της αρεσκείας σου, χρωμάτισε τις τρεις γωνίες του με διαφορετικά χρώματα και σκίσε τες. Βάλε τες δίπλα δίπλα με τις κορυφές στο ίδιο σημείο. Τι σχηματίζουν; Δοκίμασε και με ένα αμβλυγώνιο τρίγωνο.

\ex Σχεδίασε τρίγωνο $\mathrm{ΑΒΓ}$ με $\mathrm{ΒΓ}=6$ cm, $\hat{\mathrm{Β}}=50^\circ$ και $\hat{\mathrm{Γ}}=60^\circ$. Υπολόγισε τη γωνία $\hat{\mathrm{Α}}$ και μετά μέτρησέ τη. Συμφωνούν;
